\section{Contexto social}
% Empezar con el contexto social, económico, global del problema con sus debidas referencias

En el marco de la aceleración industrial contemporánea, donde la automatización y la 
robótica colaborativa se erigen como pilares de la economía 
global \parencite{ElFinanciero_Renacimiento_2025}, la consolidación de la Industria 4.0 
exige una transformación educativa paralela (denominada Educación 4.0) para dotar al 
capital humano de competencias tecnológicas y críticas. Este esfuerzo 
ha llevado a instituciones líderes en México a adaptar 
sus planes de estudio a las necesidades de la 
Industria 4.0 \parencite{Avalos_IPN_2023, Zazueta_Educacion40_2023}. 
Esta transformación no solo impulsa la eficiencia en la manufactura, sino que, alineada con las tendencias 
de la robótica industrial hacia 2026 \parencite{MMS_Tendencias_2023}, convierte la capacidad 
de programar manipuladores autónomos en un requisito indispensable. Sin embargo, el 
desarrollo de estas tecnologías en el ámbito educativo suele toparse con una desconexión 
significativa: la brecha entre la abstracción matemática, 
indispensable para modelar y generar soluciones propias, 
y su implementación práctica en la ingeniería de mecatrónica.

En el contexto nacional, esta brecha tiene repercusiones profundas 
y arraigadas. Diagnósticos 
estratégicos que datan de hace casi dos décadas \parencite{SE_Mecatronica_2007} ya advertían de un 
desarrollo tecnológico incipiente, donde los esfuerzos se limitaban a la construcción 
de prototipos académicos mientras la industria dependía de tecnología importada. Esta 
dinámica estructural ha derivado en la actual crisis de talento digital que enfrenta la 
manufactura \parencite{Energy_Talento_2025}; análisis recientes de política industrial señalan que, 
pese al potencial exportador del país, apenas el 7\% de la fuerza laboral posee capacidades 
digitales avanzadas \parencite{Solleiro2023}. La escasez de mano de obra calificada para 
diseñar y controlar sistemas de automatización avanzada no solo limita el campo 
profesional a tareas de mantenimiento y soporte, sino que pone en riesgo de 
desplazamiento laboral a sectores de la población que carecen de estas competencias 
de alto nivel \parencite{Zazueta_Educacion40_2023}. Ante la urgencia de revertir esta 
tendencia y fomentar el desarrollo de tecnología propia, las instituciones educativas 
deben buscar estrategias para generar herramientas de software avanzadas que cierren 
esta brecha. Este proyecto terminal surge precisamente como una respuesta a dicha 
necesidad, materializándose en una iniciativa de vinculación académica entre 
la Universidad Abierta y a Distancia de México (UnADM) y 
el Tecnológico de Estudios Superiores de Huixquilucan (TESH),
cuyo objetivo es aplicar la solidez conceptual en matemáticas de la 
UnADM para la creación conjunta de algoritmos de planificación de 
trayectorias que transformarán a los ingenieros del TESH de 
operadores de sistemas a desarrolladores de tecnología.

Este esfuerzo por integrar la conceptualización matemática y
la visualización computacional se enmarca, sin embargo, 
en un desafío estructural adicional: la limitación de recursos materiales.
Dado que la histórica falta de 
financiamiento ha restringido el desarrollo de la robótica en 
México \parencite{UV_Recursos_2016}, una situación agravada por una 
estrategia digital nacional que prioriza la austeridad sobre la inversión en 
modernización industrial \parencite{Solleiro2023}, 
y que la complejidad para adquirir infraestructura especializada en instituciones públicas 
limita la creación de laboratorios 4.0 \parencite{Avalos_IPN_2023}, el TESH 
ha adoptado una filosofía de resiliencia tecnológica basada en el desarrollo propio de herramientas de simulación 
y hardware de bajo costo \parencite{Garcia_Simulacion_2025}. 

Bajo la guía del Dr. Enrique García Trinidad y el departamento de Ingeniería 
Mecatrónica, la comunidad estudiantil diseña y construye sus propios manipuladores y entornos virtuales. 
Esta estrategia no solo mitiga la carencia económica, sino que la transforma en una 
fortaleza pedagógica. Al obligar a los alumnos a comprender el funcionamiento 
electromecánico de los robots desde sus cimientos, en lugar de limitarse a operar cajas negras 
importadas, el TESH exige un dominio conceptual que es esencial para la creación de tecnología propia.

En este esquema de cooperación, la UnADM complementa el esfuerzo de 
hardware de bajo costo del TESH con la fundamentación teórico-matemática 
necesaria. Esta sinergia se manifiesta en que la formación en matemáticas 
avanzadas (esencial para el modelado de espacios de configuración y 
estructuras toroidales) habilita la capacidad de abordar 
la planificación de trayectorias no solo como un problema de control, 
sino como una aplicación directa de análisis de conectividad sobre 
plataformas físicas reales. De esta manera, el proyecto desarrolla el 
pensamiento crítico y la capacidad de resolución de problemas complejos, 
capacidades esenciales de la educación 4.0 \parencite{Zazueta_Educacion40_2023}. 
El resultado de esta colaboración es una contribución social y académica 
doble: se valida la aplicabilidad de la conceptualización matemática 
en la ingeniería y, se dota a los estudiantes del 
TESH de herramientas de software avanzadas (Python) y 
algoritmos propios. Esto los posiciona no como técnicos de mantenimiento, 
sino como desarrolladores de tecnología capaces de enfrentar los 
retos de la Industria 4.0 y mitigar el riesgo de obsolescencia laboral.

\section{Planteamiento del problema}
% Enfatizar el problema de la organización desde el punto de vista matemático

Matemáticamente, el movimiento de un manipulador antropomórfico como el 
Brazo Robótico Comunitario (\textit{Community Robot Arm}) \parencite{community_robot_arm_20sffactory} se modela como la evolución de un punto dentro de un 
espacio de configuración con estructura toroidal \parencite{Munkres2000}. 
Esta geometría no es arbitraria, sino que surge de la naturaleza rotatoria 
de sus articulaciones: el estado del sistema no se desplaza sobre un plano 
euclidiano infinito, sino en un entorno cíclico donde el límite de rotación 
coincide con el origen. El desafío fundamental radica en dotar de estructura 
a este espacio para garantizar la conectividad 
por caminos \parencite{Munkres2000}, transformando así la geometría física 
del robot y sus obstáculos en una entidad matemática navegable.

La dificultad para materializar esta transformación matemática 
es particularmente notoria en el entorno académico del TESH, 
donde faltan plataformas que permitan analizar programáticamente 
la interacción entre la geometría del espacio de 
configuración y los métodos de búsqueda de trayectorias. Esta dicotomía 
se hace evidente al contrastar dos filosofías de exploración: el 
algoritmo A* \parencite{Hart1968}, que impone una organización en 
retícula determinista mediante una partición discreta en 
voxeles \parencite{Choset2005}, y el algoritmo RRT \parencite{LaValle1998}, 
que utiliza un muestreo estocástico para construir árboles conexos. 
Por consiguiente, el reto es transformar la estructura toroidal de un 
concepto abstracto a una herramienta de cálculo numérico, permitiendo 
estudiar cómo la discretización del espacio y el azar del muestreo determinan 
la viabilidad de una trayectoria segura.

A pesar de la existencia teórica de estos algoritmos, su implementación 
en manipuladores específicos presenta una barrera operativa crítica: 
la ausencia de herramientas que permitan a los investigadores observar 
y analizar el comportamiento de estos procesos dentro del espacio de 
configuración. Actualmente, la experimentación suele limitarse al uso de 
frameworks de alto nivel como MoveIt en ROS2 \parencite{macenski2022robot} y su visualizador nativo RViz2 \parencite{rviz2_software}. Si bien estas herramientas 
son estándares industriales, actúan como \enquote{cajas negras} que resuelven la 
trayectoria automáticamente, ocultando la complejidad matemática subyacente 
e impidiendo que el estudiante comprenda la eficiencia de la ruta o la 
viabilidad geométrica de la misma. Esta limitación didáctica inherente 
a las herramientas comerciales exige un cambio de paradigma.

El enfoque ingenieril tradicional se basa en soluciones trigonométricas 
específicas para cada modelo de hardware; en contraste, la transición 
hacia una fundamentación matemática avanzada permite generalizar el 
problema de la movilidad, tratando al manipulador como una variedad 
topológica independiente de su arquitectura física. Este cambio de paradigma
es necesario para trascender el cálculo cinemático convencional, 
permitiendo que el software analice la estructura global del espacio libre 
y garantice la convergencia de los algoritmos en escenarios de alta 
dimensionalidad. Para validar esta abstracción, se utiliza como prueba 
y aplicación directa el Brazo Robótico Comunitario de hardware abierto, el cual 
sirve como el escenario ideal para mapear la teoría de conjuntos y la 
topología en un entorno de cálculo real y de bajo costo.

Por lo tanto, el problema central no radica únicamente en la movilidad 
del hardware, sino en la carencia de un entorno de software que permita 
la experimentación y el análisis algorítmico, fundamentado en un modelo 
matemático riguroso, para la toma de decisiones en la planificación de 
trayectorias sobre plataformas físicas reales.



\section{Justificación}
% Describir el impacto del proyecto a la resolución de la problemática y sustentar el uso de los modelos matemáticos a usar.
La relevancia de este proyecto reside en su capacidad para convertir 
la gestión de trayectorias robóticas en un entorno de experimentación 
geométrica rigurosa. Desde una perspectiva teórica, el modelo 
se cimienta en la formalización del espacio de configuración como un 
espacio métrico con estructura toroidal, lo que permite tratar el 
movimiento del brazo como un problema de planificación de movimiento 
topológico \parencite{farber2008}, donde la búsqueda del camino óptimo 
está intrínsecamente ligada a las propiedades de la variedad.

Asimismo, la aplicación de la teoría de recubrimientos proporciona un 
sustento matemático robusto para la detección de colisiones. Esto se logra 
mediante la aproximación del volumen del manipulador con uniones finitas de 
bolas abiertas \parencite{coumar2025foam}. Este enfoque asegura que 
la discretización del camino sea topológicamente consistente con una 
trayectoria continua en la variedad, utilizando el estudio de las 
inestabilidades topológicas en el mapa cinemático \parencite{Mosquera2017} 
para asegurar la estabilidad del sistema y la evasión de configuraciones 
críticas.

De este modo, el proyecto fundamenta cada decisión algorítmica en 
principios de geometría computacional y topología algorítmica, 
ofreciendo una plataforma en Python que valida la fidelidad y eficiencia 
de la búsqueda en espacios no euclidianos. Esta síntesis entre la 
abstracción conceptual y la herramienta de cálculo no solo resuelve 
la problemática de software, sino que eleva la capacidad de los futuros 
ingenieros del TESH para innovar en la robótica nacional mediante el 
uso de algoritmos y modelos matemáticos propios.




\section{Pregunta de investigación} 
% (1 o varias)

¿De qué manera es posible modelar el espacio de configuración de un manipulador robótico de $n$ grados de 
libertad mediante una caracterización de variedad toroidal para garantizar la conectividad por caminos en un subespacio 
libre de obstáculos, integrando la teoría de recubrimientos por bolas abiertas y el análisis del mapa cinemático en 
una plataforma de software que permita el estudio comparativo de la eficiencia entre los algoritmos A* y RRT?

De esta pregunta central se desprenden las siguientes interrogantes que guían el análisis matemático del proyecto: 
\begin{itemize}
 
    \item ¿En qué medida la selección de una métrica específica ($L_1, L_2$) o la resolución de la discretización 
    en el espacio de estados condiciona la convergencia hacia la trayectoria óptima en el algoritmo A*?
 
 \item ¿Cómo influye el muestreo estocástico del algoritmo RRT en la completitud probabilística de la búsqueda al 
 navegar una geometría toroidal con restricciones de colisión complejas?
 
 \item ¿De qué manera la aproximación del manipulador mediante la unión finita de bolas abiertas optimiza el 
 costo computacional de la detección de colisiones sin comprometer la consistencia topológica del modelo?
 
 \item ¿Cómo facilita el desarrollo de una herramienta de visualización propia en Python la 
 transición de los estudiantes del TESH de una comprensión operativa (caja negra) hacia un análisis crítico de los 
 fundamentos matemáticos de la robótica?
\end{itemize}

\subsection{Hipótesis} 
%(1 o varias)

Se propone que la formalización del espacio de configuración como una variedad toroidal, sujeta a una discretización métrica 
y al modelado mediante recubrimientos de bolas abiertas, permitirá el desarrollo de una 
plataforma en Python capaz de resolver la conectividad por caminos en subespacios libres de obstáculos. 
En este contexto, se anticipa que el algoritmo A* garantizará el hallazgo de la trayectoria óptima respecto a la métrica 
seleccionada ($L_1$ o $L_2$) dentro de la discretización del grafo, siendo ideal para análisis de precisión, 
mientras que el algoritmo RRT ofrecerá una completitud 
probabilística y una mayor eficiencia computacional en entornos donde la partición en retícula 
resulte ineficiente debido a la complejidad geométrica o la dimensionalidad del espacio.. 
El uso de estas abstracciones, implementadas bajo un diseño de software modular y orientado a la robustez lógica, proporcionará una herramienta 
de análisis que minimice la incertidumbre en la detección de colisiones y permita la evasión de inestabilidades 
topológicas mediante el estudio del mapa cinemático, optimizando así la capacidad analítica y 
el aprendizaje de los algoritmos de búsqueda para los estudiantes del Tecnológico de Estudios Superiores de Huixquilucan.

\section{Objetivos} 
% (usando verbos en infinitivo)
\subsection{Objetivo general}
% Contestar las preguntas ¿qué?, ¿cómo? y ¿para qué?

Desarrollar una plataforma de planificación de trayectorias para el manipulador 
Community Robot Arm, basada en el modelado topológico de su espacio de configuración, 
con el fin de comparar la eficiencia de algoritmos de búsqueda y fortalecer la enseñanza de 
la matemática aplicada en la ingeniería.

\subsection{Objetivos específicos}
% Lista de objetivos para lograr el objetivo general

\begin{itemize}
\item Caracterizar la estructura toroidal del espacio de configuración ($T^n$) y 
definir analíticamente el subespacio libre $\mathcal{C}_{free}$ mediante la teoría de 
recubrimientos, utilizando uniones finitas de bolas abiertas para optimizar la detección 
de colisiones.

\item Implementar los algoritmos A* y RRT mediante un paquete personalizado 
en ROS2, integrando el análisis diferencial del mapa cinemático para garantizar 
trayectorias estables y libres de singularidades topológicas.

\item Evaluar comparativamente el desempeño de ambos algoritmos en términos de 
convergencia y optimización métrica ($L_1, L_2$), diseñando un sistema de 
visualización en RViz2 que traduzca procesos abstractos en herramientas 
didácticas de experimentación.

\end{itemize}
